\section{Explicit public development objectives}
Twelve newly authored public examples make three original development objectives
concrete using their retained public titles and prefix-source APIs. They are not
recovered historical issue text, hidden acceptance checks or independent families.
On unchanged XML, OAuth and Click prefixes, respectively, 2/4, 1/4 and 2/4
examples conform. XML attaches namespace declarations to a descendant when the
declaring element has no ordinary attributes; OAuth authorization metadata raises
an AttributeError, including with explicitly supplied response types; Click help
still enumerates choices when their display is disabled. Public comparison cases
exercise ordinary attributes, disabled namespace processing, token-only metadata,
visible choices and explicit metavar behavior.

All twelve examples run in three fresh processes under frozen 60-second bounds,
with no setup errors. An earlier run's three missing-constructor-argument errors
are retained as harness failures. The revised OAuth setup uses actual token and
endpoint classes and all required combined-server claims. Full unchanged source
inventories retain 16, 199 and 145 entries; actual imported modules bind to those
prefixes. A separate 545-check audit reconstructs predicates and source/process
provenance without further execution. The new objective manifest binds text,
examples, repair surfaces and public API sources. Together with the separately
retained nine-example Pyparsing public-report reproduction, these provide explicit
development objectives for the original four-family matrix. Independent scorer
admission, complete operational trajectories and matched cache/component arms
remain required. No model or hidden scorer is called; example conformity is not
a repair-success or system-quality estimate.

A later qualification exposes a limit in the newly authored OAuth expectation.
The original accepted reference passes three of four public examples; combined
server metadata omits one configured response type. An independently authored
configured-key-preserving candidate passes all four new examples but fails the
unchanged original acceptance check. Public-only API review does not justify
requiring the abstract-validator metadata probe to advertise every internal
dispatch key. Exact three-key equality is therefore an unqualified author
assumption, retained diagnostically rather than imposed on original acceptance.
This is a prospective contract correction before new model dispatch; neither
the original criterion nor the contrary observations are overwritten. It is not
a proof that every possible repair faces incompatible requirements. The
supplemental native call's missing enclosing timer also remains unknown.

\section{Additional benchmark qualification}
\label{sec:incremental-recruitment}
An independently pinned SWE-bench Verified metadata projection~\citep{swebench2024,swebenchverified2024} supplies 500
candidate repairs from twelve repository names. Ten of eleven prospectively
selected two-issue pairs have forward base-commit ancestry; one has reverse
ancestry in PR-date order. A separate current GitHub identity audit checks all
192 candidate/original repository pairs, finds no shared repository or declared
fork identities, and flags two shared-owner groups. This does not establish
historical code independence, portfolio separation or absence from model
training. A later audit reports test-specification problems in a selected hard
subset and evidence of solution exposure in tested frontier models~\citep{swebenchverifiedaudit2026}.
Its selected-subset rate is not a whole-dataset error estimate, and it does not
measure contamination in our provider profile. Public examples exposed during
this literature review are separately recorded before any prospective split.
No new confirmatory family or complete natural trajectory is enrolled.

Requests 6028 has a separately named ARM environment preserving its
public Python version and thirteen listed Python package versions. The unchanged
native SWE-bench 5.0.2 parser/grader observes all 187 designated checks in both
controls: unfixed source passes 185 and fails the two required failing checks;
reference source passes all 187 and is the only resolved control. Both collected
outputs also contain eight missing-fixture setup errors outside the designated
set. The original script lacks a test-exit marker, so a zero container exit
cannot establish acceptance. This is designated-check qualification, not an
error-free full-suite result or equivalence to the published x86 environment.
Protected reference/test content remains in curator custody; model-facing
selection reads public metadata only. Earlier import and decoding failures are
retained. Infrastructure qualification contributes no provider-token or repair
success observation for either treatment.

Requests 5414 subsequently qualifies in a new named ARM profile
using its public development dependencies and Docker's default bridge network.
Both controls complete all 131 designated checks: unfixed source passes 130
and the reference passes all 131. The full parsed outputs have 288 versus 289
passes, respectively, with one expected failure each and no setup errors.
The original offline port, a dependency-build failure and a plugin-startup
failure remain retained; the successful profile is a development correction.
This does not make the original failed port successful. A later candidate
grading service version fails its four-control qualification: three controls
behave as expected, but this node's reference is not accepted despite all
designated checks completing. Four existing public network-timeout checks
instead encounter a network-unreachable condition; the earlier standalone
run used the same image, package versions and evaluation bytes. The unique
cause and a reliable remedy are not established by this discrepancy.
Eight service rejection controls pass. This retained discrepancy illustrates
why standalone benchmark qualification alone does not qualify a new endpoint.
A subsequently frozen service version passes one attempt at each of its four
native controls and all ten rejection controls. It preserves exact objective,
profile, cumulative patch and resulting source-tree bindings. Public network
readiness probes bracket Requests 5414's evaluation, with no automatic retries;
an incomplete grade remains unavailable for repair decisions. All eight
qualification probes pass. The 94-check audit retains 1,175 artifacts; the four
control containers take 87.429894 seconds in total. This qualifies the named
service and its observed controls, without proving the unique cause or remedy
of the previous network failure or stability throughout future test runs.
Candidate and benchmark tests still share the original test process; adversarial
oracle integrity is not established. Requests 6028 retains its eight
outside-set fixture errors.

An additive public-only cumulative-source control changes two production files
and adds one Git-ignored Python file. The service inventories all 98 resulting
files, matches the independently expected tree, and returns the unchanged
unfixed grade of 185/187 in 6.048 seconds. This checks complete source-carry
transport into grading, not successful repair or an accepted multi-issue
trajectory. Earlier four-control cohorts remain unchanged.

A further pinned metadata-only projection yields 21,336 SWE-rebench
rows~\citep{swerebenchmetadata2026} from 3,468 repository names and 731
SWE-bench Pro rows~\citep{swebenchprometadata2026} from eleven names.
The 22,067-row sum is not deduplicated across datasets; protected issue,
reference and test fields are not decoded in this projection. A prospectively
seeded screen uses 2025 metadata dates, available image fields, distinct base
commits and exclusion of prior repository/owner names. It yields 53 candidate
pairs. All first sixteen pairs have actual forward Git ancestry, and 448
comparisons against the 28 prior current repository identities find no
repository, declared-fork or owner match. Each candidate has a distinct current
repository and owner identity. These checks do not prove historical code
independence, portfolio separation or training novelty. Image-field eligibility
is a convenience filter; all 32 selected image names specify x86\_64, whereas
the evaluation host is ARM. Source/environment/scorer qualification and a
prospective confirmatory split remain open; zero new families are enrolled.

The first frozen additional pair, Zarr 2996 and 3015, subsequently passes four
original native grading controls on explicitly ported ARM environments. Unfixed
source passes 16/22 and 157/158 designated checks, respectively; each reference
passes its complete set. Every control runs once, with zero retries or provider
calls. The native test-spec factory, test commands, parser and designated
criteria remain unchanged; the script port removes four Conda activation
commands. The environments preserve Python 3.11.11/3.11.14 and all 50/67
observed installed distribution versions, while architecture, platform binaries
and system-library builds differ. The local native fork reports version 4.0.3.
Its complete declared dependency installation failed because an ARM YSON
binary distribution was unavailable; a separately named local Docker backend
uses the available Python implementation, leaving cloud backends unsupported.

The four native grader container intervals total 10.647 seconds, excluding
environment acquisition, builds, Git preparation and audit. Earlier missing
build-helper, project-version and Git-image-reference failures remain retained.
The full public collection probe collected 3,045 tests for 2996 and exited zero;
3015 collected 3,123 but exited one during teardown with resource/unraisable
warnings. The targeted controls do not reproduce that teardown failure, and
their success does not erase it. As with Requests, the original script lacks
a separate test-exit marker, and native shared-process execution does not
establish adversarial oracle integrity. These results qualify the four original
controls; candidate endpoint, complete trajectory and token economics require
separate evaluation. No confirmatory family is enrolled by these controls.

A subsequent Zarr candidate endpoint qualifies all nineteen frozen controls,
with 190 audit checks: six actual native grading calls, ten rejections before
dispatch, two wrong-source containers stopped before hidden grading, and one
duplicate-ID rejection. Empty and cumulative-public candidates retain the
unfixed grades; reference production projections pass 22/22 and 158/158.
Each projection preserves the full original reference's production and
base-public bytes while omitting one ancillary added document, then reruns
the unchanged designated checks. This qualifies the projected candidates,
not transport of the full original patches. Non-Python edits remain outside
this endpoint's admitted scope.

Exact public base inventories contain 189/181 files, including the editable
installation's generated version module. Cumulative controls make three
production edits, including a Git-ignored addition, and retain all 190/182
resulting files. The initial endpoint attempt failed before hidden grading
because of raw mode differences and its rejection of the ancillary document;
a later preparation attempt caught a no-op Git apply in a nested checkout.
Both failures remain retained. The qualified revision preserves image-exact
modes and verifies actual edits in isolated Git repositories.
Six grader containers total 14.426 seconds within 30.363 seconds of outer
service calls, including the duplicate-ID control. Preparation, environment
acquisition and audit are separate. There are zero retries or provider calls.
Context, accepted-source trajectory carry, publication and model efficacy are
not enrolled by this endpoint qualification; adversarial oracle integrity
and non-Python transport remain unqualified.

The second frozen additional family, Xradar, then completes its four original
native controls on explicit ARM ports. Issue267 has 11/13 unfixed passes and
13/13 reference passes; issue269 has 55/58 and 58/58. All designated cases
complete, with one attempt per control and no provider calls. Full original
reference patches, including ancillary edits, are retained. The unchanged
native factory/parser/test commands govern both nodes. Public collection yields
354 and 356 cases with exit zero. Each public prefix matches 106 original Git
blobs, and all 22 data fixtures selected from pristine public tests match their
registry hashes, totaling 183,051,686 bytes; both images fetch them offline.

Both ports preserve Python3.10.17 and all 196 installed distribution records,
including two project distributions with natural imports resolving to the
base-specific editable source. Four exact package versions require ARM source
wheels. Conda binary/build variants and optional compiled-feature parity remain
unqualified. Retained failures include initial duplicate-project metadata
ambiguity, a missing ARM binary wheel, and a node-orchestrator key error after
the shared runtime built. Versioned corrections preserve runtime pins and do
not retry grading. Failed and successful shared builds take 118.916 and 188.913
seconds; cached-wheel materialization takes 19.143 seconds, and the complete
node-build/import/cache/package/collection phase 31.950 seconds. Four native
grader containers total 69.044 seconds; these are scoped intervals, not complete
recruitment cost. Environment and native audits pass 373 and 129 checks,
respectively, with 120 custody artifacts and 59 native source files. Candidate
endpoint, context, model trajectory, confirmatory enrollment and adversarial
oracle integrity remain unqualified.

For a later development experiment on the existing pyparsing snapshot, the
current public issue report supplies a concrete objective. Its nine examples
reproduce the reported recursive-whitespace defect on unchanged public source:
eight conform to the reported expectation and one fails. The report is frozen
for that new experiment; its exact bytes at the original base date remain
unverified. A reproduced public example is not complete hidden acceptance.

\section{Repository preparation and source custody}
\label{sec:incremental-context-preparation}
\paragraph{Repository context custody and preparation cost.}
A separate native producer/consumer qualification retains all 302 regular files
across the XML and pyparsing development snapshots, including unparsed source.
Its bounded semantic scope analyzes one of six and five of 154 Python files,
respectively; four dependency capsules grant no source omission. Four
cold/incremental root-and-block comparisons agree, four exact source expansions
execute, and 22 negative controls reject. An added caller checks required-source
CIDs and current bytes because native freshness omits that mapping; native
``sufficient'' flags are not behavioral authority. Full pyparsing scanning
produces 3,957 symbols and 55,859 edges, but the unmodified producer process does
not finish within its 600-second bound. Hoisting a repeated immutable
node-CID index preserves the full XML state/root and reduces its preparation
from 64.838 to 26.296 seconds in one observation per profile. Full pyparsing
still reaches the 600-second bound after that change. Source retention and
bounded integration therefore do not establish affordable whole-repository
preparation, context compression or repair efficacy. Three later isolated
profiles add a fixed-profile CID encoder, cached syntax validation and an exact
literal-union predicate. Across all five profiles XML state/root identity is
unchanged and every native consumer passes; each full pyparsing producer still
hits the process deadline. In the last profile the producer writes its result
at 597.222 seconds before the 600.202-second termination; the native consumer
does not run. An earlier stage diagnostic had also assembled and verified a
full bundle before timing out in source-binding reads. Missing final output
therefore does not prove that bundle construction itself failed.

The literal-union change preserves all 55,859 actual link predicates against
3,957 node CIDs, an embedded-CID rejection, 1,235 constructed comparisons and
64 native regressions. The independently validated full bundle has 67,872
blocks and agrees with the last producer's state/root. Its 72.358 versus 2.834
seconds measure only one diagnostic predicate pass per implementation on a
shared host; they are not a whole-pipeline speedup estimate.

A generic current-repository service separately checks exact source inventory,
file modes and objective/request identities. Twelve controls pass across three
XML preparations: a source change requires new state/root/pack identities,
whereas an objective-only change retains state/root and creates a new pack.
Both public Requests nodes then complete actual native producer/consumer
roundtrips, retaining 97 files each and scanning all 35 declared Python files.
Eleven negative controls pass per node; neither node permits any raw-source
omission. These are context and custody checks, with zero provider calls and
no repair-quality or compression comparison.
Four subsequent Requests API expansions return 79,942 exact raw bytes across
two real read processes, preserving source, request, freeze and objective
identities. They add no omissions or model calls.

A sixth isolated profile memoizes identity only within an exact immutable
native state instance under frozen codec/runtime bindings. Both full producer
processes and their separate native consumers now complete. XML preparation
takes 9.103 seconds; pyparsing takes 566.254 seconds, with
553.516 CPU seconds. The latter producer process takes 569.609
seconds and its consumer 1.740 seconds. All 302 regular files and
11,781,935 bytes are materialized; pyparsing's additional opaque symlink is
retained without following it. Both sources' previous available state/root
identities agree, including the independently checked 67,872-block pyparsing
bundle. All 22 binding/tampering controls pass and all 7,255 frozen runtime
source files remain unchanged. Zero source omissions and zero signed receipts
are retained, along with the unresolved semantic frontier.

The profile passes 159 native regressions. Separate full XML unchanged,
comment-change cold and comment-change incremental states preserve serialized
identities and all 1,992 reachable blocks across the parent and memoized
profiles. Each profile performs twenty actual source materializations and
rejects twenty corrupt-source reads. This qualifies immutable-object identity
caching, not execution/proof reuse or a defense against reflective mutation of
frozen objects. The successful full run establishes completion within the bound, not a
replicated speedup or token benefit; the five previous deadlines remain results.

A subsequent frozen full-pyparsing comment-change comparison fails cold versus
incremental equality. Its three workers complete within a separately declared
900-second bound: baseline 656.008 seconds, changed cold 609.801 seconds and
changed incremental 784.021 seconds. Baseline and changed-cold workers run
concurrently on the shared host; incremental restoration, full source guards,
native previous-state/bundle loading, persistence and process exit are included
in its worker interval. Earlier driver enrollment copy/freeze creation is not
separately timed. These observations establish no incremental speedup.

The changed snapshot, symbols and artifacts agree, but 1,568 of 55,859 edge
records differ: 1,564 ambiguous candidate edges change resolution metadata from
\texttt{finite\_may} to \texttt{definite}, and four derived references change
with their edge identities. All retain conservative confidence and multiple
candidates. Consequently state and semantic-root identities differ. Independent
comparison reads every persisted block: 65,856 shared CIDs have identical bytes,
with 2,016 distinct CIDs on each side, covering links, capsules and their indexes
and roots. Both stores contain 67,872 blocks; their payloads total 112,728,034
cold versus 112,724,906 incremental bytes. Ten source-read controls pass and
7,255 runtime files remain unchanged. The driver exits one at the equality
assertion after all workers exit zero; this is a retained graph-metadata failure,
not another preparation timeout. A later actual native explanation-consumer
probe loses the ambiguity limitation on all 1,564 directly changed edges while
retaining their conservative-confidence limitation; four derived references
are excluded. The fresh process takes 3.070 seconds and rechecks 7,255 runtime
files. This establishes an explanation-output consequence, not an execution
authorization or soundness failure.

A subsequent isolated guard refuses previous-state frontend reuse when graph
projection has already replaced a lexical target with a concrete symbol.
Such paths are reanalyzed from current source before whole-inventory resolution;
the guard does not simply preserve stale resolution metadata. Seven changing
inventory controls preserve state, roots, all block bytes and explanations.
The full pyparsing comment-change comparison then restores exact serialized
state/root equality and all 67,872 block payloads (112,728,034 bytes), also
matching the earlier cold output. It reuses 20 of 154 Python paths, 172 symbols
and 1,428 original frontend edges, while rescanning 134 paths. These counts
are not omitted source bytes or avoided execution checks.

The two fresh workers run concurrently under unchanged 900-second bounds:
cold takes 437.699 seconds and guarded incremental 614.522 seconds. The latter
includes 107.173 seconds of prior-state restore/reverification and 55.736
seconds of scanning; graph construction and full verification still execute.
The earlier 656.008-second baseline preparation is reused as an artifact and
remains a separately incurred cost. Enrollment copying/freezing is not timed;
the subsequent independent byte audit takes 3.625 seconds. Nine new source
controls pass and 7,255 runtime file digests remain unchanged. Native regression
tests retain 187 passes and one failure: an existing performance expectation
requires reuse of a path that this correctness guard deliberately rescans.
An earlier invalid authored fixture is also retained. A further inventory
control exposes the guard's limit: adding a previously missing imported module
leaves incremental confidence conservative while cold analysis becomes exact.
The guard's successful comment-only comparison remains valid, but it is not a
general lossless solution. No execution-soundness consequence is inferred.

A new profile instead persists the original Python frontend observations before
pytest unification and graph projection. It binds their exact sidecar bytes to
the accepted parent state, source inventory, namespace and extractor/runtime
profile; unsupported or stale inputs fall back to current analysis. Current
pytest analysis and whole-inventory graph resolution still run. The full
comment-change comparison reuses 153 of 154 Python paths, 3,885 original symbols
and 52,977 original edges, extracting 772 edges from the changed file. Cold and
incremental serialized state, root, all 67,872 blocks/112,728,034 bytes and the
complete sidecar agree. The native state and blocks also equal the earlier cold
output. Each sidecar adds 53,753,923 bytes of storage. Ten bounded equality
cases and sixteen rejection/fallback checks include the earlier unresolved-import
confidence case; ten full-worker source controls pass and 7,256 runtime hashes
remain unchanged. Accepted honest producer custody remains an assumption:
hash binding alone does not prove that arbitrary observations follow from source.

All three new workers finish within frozen 900-second bounds. Baseline sidecar
bootstrap takes 231.354 seconds in addition to the separately paid original
656.008-second preparation. Changed cold takes 440.276 seconds and incremental
607.971 seconds. Incremental includes 105.262 seconds of parent restoration and
verification, 55.520 seconds of scanning with observation decode/encode,
390.254 seconds of bundle construction, 48.819 seconds of full verification and
4.727 seconds of persistence. Bootstrap overlaps changed cold; incremental
follows bootstrap, for 839.328 seconds of new worker schedule time. Enrollment
takes a separate 0.533 seconds and the independent audit 5.958 seconds.
Nested stages are not added again to process totals. This qualifies actual
original-observation reuse on one full public change; incremental preparation
is still slower, and zero source omissions or provider calls establish no
token saving or general speedup.


\section{Combined controller, publication and recovery}
\label{sec:incremental-controller-publication}
\paragraph{First combined controller qualification.}
A new caller joins actual native context reads, source-bound semantic analysis,
symbolic routing, compiler-mediated world memory, isolated behavioral scoring
and a durable effect journal. On the four existing constructed routing tasks,
two supported initial repairs pass and two remain residual-pending with the
provider disabled. Authored follow-up events consume accepted source, recheck
a comment change and reject a stale local parent. Ordinary no-op detection
receives the same current behavioral checks as the memory path.
Three actual process interruptions exercise validated continuation, durable
receipt reconciliation and an executed effect whose completion remains unknown;
the unknown effect is held without redispatch. Across thirteen events there are
nine local acceptances, two pending residuals, one stale parent and one unknown
effect. These event counts are not thirteen independent repair tasks. Thirty
scorer attempts yield 29 durable terminal receipts. Local ledger CAS is not
native Kit publication; sealing, lifecycle reuse and certified procedures are
not yet combined here. With zero provider calls this qualification establishes
no integrated token-savings effect.


\paragraph{Subsequent native publication profile.}
The separately frozen eight-event sequence yields six acceptances, one stale
branch and one pending residual. A full checkpoint at generation 0 is followed
by a delta at generation 1. A real service exits with code 94 after native
seal persistence and before CAS. A fresh process loads the pending seal,
rechecks the entire parent chain and exact local evidence, replays visible
and audit behavior, then completes native CAS/WAL through an explicitly added
development adapter, reaching generation 2. Another fresh event consumes the
native stored accepted source and publishes generation 3. A branch based on
generation 0 reaches native publication and loses with a generation mismatch.

\begin{table}[htbp]
\centering
\caption{Actual work in the eight-event native publication qualification.
Event outcomes are dependent continuation observations.}
\begin{tabular}{lr}
\toprule
Operation & Count \\
\midrule
Controller invocations / terminal effects & 9 / 68 \\
Native context reads / source expansions & 4 / 8 \\
Input scans / accepted-candidate scans & 8 / 3 \\
Symbolic attempts / produced candidates & 4 / 3 \\
Native routing calls & 8 \\
World-memory service calls / compiled actions & 11 / 2 \\
Publication attempts / native reads / recovery & 7 / 4 / 1 \\
Controller / additional admission scorer processes & 18 / 24 \\
Provider requests / provider tokens & 0 / 0 \\
\bottomrule
\end{tabular}
\end{table}

Invocation wall time is 60.906438 seconds; the additional 24 behavioral
replays consume 13.759197 seconds within that total, not as a separate additive
cost. Context enrollment, inventory freezing, implementation and read-only
audit are outside this runtime denominator. Crash-time work is retained through
a durable receipt, and recovery does not double-count its aliased completion.
Both ordinary and memory paths run the same cheap current-source check and
accept no-op revisits without symbolic or model repair. The compiled memory
action dispatches an actual check; no reduction in checks is demonstrated.

The complete required leaf set is externally frozen: accepted source,
semantic preparation, effect inventory, visible behavior and audit behavior.
Every changed candidate receives a new native scan. All five leaves are
replaced per delta. Native integrity verification checks content and digests;
unsigned local receipts additionally depend on trusted controller/scorer
execution, a durable effect ledger and an evaluation-owned CID index. This
is neither a theorem about arbitrary Python nor production trust admission.
Stored parent chains are recursively checked and the requested source is
behaviorally replayed on read. Broader crash phases, concurrent publication,
provider effects, real repository edits and lifecycle reuse remain open.

The read-only audit passes 6,424 checks. The archive retains 504 source and
metadata files, including 483 native source files across the loaded closure
and preceding context enrollment. The first profile's omitted imported wrapper
and the first audit's pointer-envelope parser failure are retained. Correcting
the audit did not rerun behavioral outcomes; the archived previous controller
still passes its original 9,188 checks.

\paragraph{Full-repository service fixture and measured cost.}
A separate curator-held Requests 6028 fixture begins with its reference
candidate and then adds a neutral comment to the same accepted production
source. Both states contain all 97 source files. It executes two full native
context preparations, four world-memory service calls and six original native
grader invocations: two initial checks and four fresh publication, interruption,
recovery and read replays. All six complete the 187 designated checks.
Reference and hidden test payloads remain in curator custody. This fixture
does not execute the Requests 5414-to-6028 upstream transition or any model
repair, and it does not establish adversarial oracle integrity.

The first full checkpoint reaches generation 0. The second process constructs
a delta and actually exits with code 94 after seal persistence but before CAS.
A fresh process recovers generation 1; another fresh process reads and restores
the exact second accepted source. Wrong-tree, wrong-receipt, source-drift and
incomplete-effect controls all reject before grader replay. All four required
leaves---accepted source, semantic preparation, observed non-sealing effects
and native behavior---are replaced on transition. The effect leaf excludes
all sealing operations, including prior reads; the outer workflow ledger must
still count their actual work. Parent/source admission and fresh grading are
additional checks. Unsigned local receipts and native integrity mode establish
neither a theorem of execution nor verification reuse.

\begin{table}[htbp]
\centering
\caption{Measured disjoint sequential intervals for the Requests 6028 curator
fixture. Grader processes are nested in these phases and are not added again.}
\begin{tabular}{lrr}
\toprule
Phase & Count & Wall seconds \\
\midrule
Full native context preparation & 2 & 318.283 \\
Initial accepted-candidate checks & 2 & 16.047 \\
World-memory service calls & 4 & 5.858 \\
Publication/recovery/read and rejection controls & 8 & 37.037 \\
\midrule
Sum of measured intervals & & 377.226 \\
\bottomrule
\end{tabular}
\end{table}

The six grader containers consume 31.625 seconds within the measured
intervals. Their four replay service wrappers are also nested, not independent
extra costs. The frozen-start-to-result artifact span is 377.640 seconds;
this is not a separately instrumented whole-process timer and excludes earlier
imports and profile creation. Later code archival and read-only audit take
0.275 seconds together, outside the fixture intervals. There is one
shared-host observation, no parallel comparison and no provider calls.
The 275-check audit retains 33 public code snapshots and hashes 1,136
private artifacts. This closes the stated service controls while leaving
multi-issue model efficacy, token savings and broader crash coverage open.

\paragraph{Natural Requests sequence: intervention and failure paths.}
The first live repository cohort contains one development family, two genuine
public issue nodes and four arms, executed in frozen D/C/A/B order with
independent local cold stores. Each trajectory has at most four provider
requests, with at most two at its first issue; unused slots carry forward.
All planning, proposal, expansion and recovery purposes share that ceiling.
The common Grok 4.6 low-reasoning profile has a 2,048-output-token limit and
a 65,536-byte request limit. Source retrieval selects whole production modules
with three-file/48,000-raw-byte limits, and exact requested expansions replace
the earlier selection. Accepted first-node source is committed on its exact
public base and carried through a real three-way merge with the next base.
No reference patch enters the model or its public history.

A uses ordinary lexical retrieval and the common current-source check. B
actually prepares and reads native context, but zero omissions qualify, so its
model context falls back to the same raw source. C additionally runs native
analysis, bounded synthesis and routing; its native candidate fails before
the model repairs the first issue. D adds actual world-memory services,
accepted-source context preparation, native publication and stored-source
read. It remains partial: lifecycle reuse, certified procedures, learned
world utility and integrated theorem proof are absent. Native permitted model
routes share the same fixed provider profile, so model-tier selection cannot
explain cost differences. Each arm retains its own histories and state; the
effect journal provides durable identities for every arm.

All four initial requests have identical canonical bytes and 13,693 provider
input tokens. Each returns one independently passing first-node repair. A and
B receive 13,568 cached input tokens each, versus 512 each for C/D; local cold
state therefore does not imply equal remote cache state. Returned reasoning
and output lengths vary. Small first-node total-token differences cannot be
attributed to source compression or a witnessed avoided model call.

A/C/D's accepted edits each conflict in the next public source merge.
Their retriever appends the required conflict module to its ordinary selection,
yielding four modules, 80,552 raw source bytes and a 91,202-byte request. The
frozen 65,536-byte gate halts before dispatch. B's different initial edit
merges without a conflict and its next request fits at 53,319 bytes. Its
three subsequent candidate grades are 184/187, 185/187 and 185/187, all
unsuccessful; the original current source passed 185/187. The second node
thus fails in every arm, through different observed paths. Source-bound
halts are retained failures, not successful symbolic closures.

The complete fixed-cohort provider totals are 14,428 (A), 65,695 (B),
14,589 (C) and 14,269 (D), summing to 108,981 across seven actual posts.
Each arm has one accepted node, so tokens per accepted node equal those
fixed-cohort totals; tokens per completed trajectory are undefined. All
unsuccessful model attempts remain charged. The first-horizon totals are
14,428/14,192/14,589/14,269 in A/B/C/D order. These are observed cumulative
costs, not counterfactual quality-versus-budget curves or independent repeats.
Charges, cached input, visible output and reasoning components remain separate
in the retained provider receipts; account settlement is not inferred.

The original driver terminates with exit 143 while B's first controller is
still running. That same worker finishes and commits its accepted event.
A separately frozen reconciliation checks the exact terminal output, journal,
source/profile identities and absence of pending effects, then executes only
the original remaining B6028 event. B5414 is never restarted. Its outer
invocation time, waited CPU and exit status remain unavailable; the controller's
own 209.770-second timer survives. B's two controller intervals sum to
408.510 seconds, which is a different timing scope from other arms' enclosing
invocation intervals. Docker workload CPU is unavailable throughout. Preparation,
grading, memory and publication/read phases are nested within runtime and must
not be added twice. Development qualification and source/environment acquisition
remain separately scoped overhead. A/C/D enclosing intervals are 96.745,
285.790 and 507.358 seconds, respectively; B's enclosing total is unavailable.

The separate ledger audit passes 1,769 source, journal, request, attribution and
budget checks. It joins all seven model receipts to their actual effects and
dispatch intents, with no unknown provider effects. Sixteen unique native
grader receipts include the two additional D publication/read replays; repeated
appearances of saved completions are not new calls. Model and effect accounting
are complete for this cohort, while outer runtime accounting remains incomplete
because B5414's enclosing timer is missing. The source archive contains 504
files. Preparation records include observed native imports; the source-expansion
process lacks a separate actual import receipt, which limits the demonstrated
execution closure. The quantitative result is not a production trust admission.

One family with implementation failures and uncontrolled server caching cannot
identify a general mechanism effect. Equal accepted-node counts also conceal
the different second-node opportunities. No arm completes the intended natural
trajectory, and the broader cache, warm-horizon, ablation and independent-family
program remains required. Subsequent transport corrections receive a new profile;
these outcomes and request bounds remain unchanged.

\paragraph{Second natural Requests cohort.}
A separately frozen conventional retrieval correction gives mandatory conflict
modules priority, followed by optional lexical source that fits the existing
three-file/48,000-raw-byte/65,536-request-byte bounds. A detached driver retains
all eight outer invocation timers without restart. On actual B/D continuations,
the conflict and two optional modules fit in 47,424 raw bytes and a 54,773-byte
request. A/C's different accepted first edits merge cleanly. All four first
requests remain byte-identical within this cohort. No source omission or
successful symbolic closure occurs; D still lacks integrated lifecycle reuse,
certified procedures and learned memory utility.

All sixteen calls execute and consume 243,255 provider tokens, including
expansion decisions, failed proposals and local transport failures. A uses
two calls per issue, while B/C/D use one then three. Each accepts only issue5414.
A's first proposal passes 123/131 designated checks; the next passes. For
issue6028, A spends one call expanding source and one on a candidate grading
185/187. C expands once and grades two candidates at 185/187. D grades its
first and third candidates at 185/187; their cumulative patch bytes are
identical. B grades one candidate at 185/187. The remaining three outputs,
one D and two B, restore all 97 pristine file contents but retain mode0644
instead of0664 on one file. The frozen transport rejects these mode-only
differences; its text-edit interface has no mode channel. Those three provider
calls remain charged, no graders run for them, and their counterfactual grades
remain unknown. No retrospective repair or rerun alters this cohort.

Fixed-cohort tokens are 56,543/62,439/63,336/60,937 in A/B/C/D order; first-node
cumulative tokens are 29,787/14,782/14,364/14,163. Tokens per completed trajectory
are undefined. Total components are 227,336 input, 1,804 visible output and
14,115 reasoning tokens; 47,744 cached input tokens are a subset. Cached-input
counts are 1,024/30,976/14,720/1,024, so local cold starts do not control provider
caching. API-reported charges total \$0.1627138000; settlement is unmeasured.

All 21 unique grader calls complete: nineteen direct and two D publication/read
replays. Six passing and fifteen failing grades include repeated validation
of the same task and are not independent success counts. Their 509.499 seconds
of container time are nested within 544.910 seconds of outer grader service,
which is itself part of runtime. The 53 unique outer effects have complete
elapsed timers. Enclosing invocations total 1,382.526 seconds. Native context
contributes 106.468/211.192/213.978 seconds for B/C/D; D publication/read,
including replays, takes 78.228 seconds. These components are not added twice.
Docker CPU and driver-preflight CPU remain unknown. The logical runtime-state
snapshot contains 5,495 files/297,870,884 bytes, excluding qualification and
private scorer stores; it is not deduplicated physical usage.

The prospectively fixed detailed audit passes 3,698 checks; independent root
aggregate and final joins reconcile all costs, nodes and 510 archived sources.
The final root join passes 1,412 checks after retaining an initial projection-ID
assumption failure and joining explicit residual/enclosing event identities.
The curator independently joins all 21 requests, dispatches, receipts and outer
completions, with no unknown effects. All 310 observed worker/world/sealing
imports match the archive; a separate native context audit binds 163 observed
producer/consumer paths and all ten context effects. Expansion lacks a separate
actual-import receipt, and standard-library/third-party import closure is not
established by these native-prefix receipts. The follow-up supports no token
saving and is not pooled as an independent confirmatory replication. Original
cache, operational-sequence and component comparisons remain required.

\section{Native proof profiles and exact statement checks}
\label{sec:incremental-native-proof}
\paragraph{A real native proof relation.}
The copied native backend was built offline with locked dependencies because
the bundled binary did not implement the source's v3 circuit. Isolated
development keys support a four-leaf Event-DAG Merkle/count relation. Three
actual proofs for counts 1, 2 and 4 verify; eight witness, public-input,
proof and key mutations reject. Primary setup takes 7.810 seconds,
proof generation 12.048--12.474 seconds, and verification
3.053--5.147 milliseconds; serialized proof JSON is 1,484 bytes.
The raw verifier accepts a different outer root that maps to the same field
element. The evaluation caller's exact-root/count and key/binary checks reject
that alternate statement before dispatch. This boundary is part of the result:
the native relation alone does not bind arbitrary outer metadata. Development
setup is not production trust admission, and a Merkle relation is not proof
that arbitrary tests executed. The v5 pytest-attestation circuit remains
unsupported in this backend.

\paragraph{The original bounded Horn profile.}
The source-built v2 backend also produces and verifies three honest implication
chains of lengths 1, 2 and 4. Setup takes 1.101 seconds, proving
2.238--2.390 seconds, and verification
3.187--5.267 milliseconds; each proof JSON is 1,700 bytes.
The honest prover rejects the tested incomplete or inconsistent inputs.
When the scalar public inputs remain unchanged, the raw verifier accepts three
changed outer metadata fields; exact statement and source/key checks at the
evaluation caller reject them before dispatch. Honest examples and wrapper
checks do not establish that every satisfying circuit witness corresponds to
the advertised textual axioms. That stronger assurance remains unconfirmed.


\paragraph{Actual accepted-history archive integration.}
A separately frozen producer/consumer profile follows the constructed
controller's native generation-3 pointer through four accepted seals, each
with five required receipt units. It excludes the stale unaccepted branch and
the separate ordinary trajectory. Producer and fresh consumer use native Kit
verified reads on exact archive clones, reconstruct the original local effect
receipts and check the complete source, parent and leaf bindings. The original
archive and controller ledger remain unchanged. No behavioral checker is rerun;
the original local execution remains the authority for its recorded observations.

Each canonical event commits its seal/parent, source and original binding,
complete receipt-CID map and verified pack digest. The already qualified
standard v3 backend proves their four-leaf archive relation. The consumer
checks the externally expected full root/count and exact binary/key, in addition
to the accepted pointer and chain. The raw scalar-root limitation still applies;
the proof does not establish that arbitrary code or tests executed correctly.
The stopped custom-v2 investigation is not used.

One honest proof and one runtime pairing verification complete. Producer and
fresh-consumer read processes take 0.325531 and 0.342170 seconds, proving
12.617511 seconds and verification 0.006725 seconds. The measured driver
totals 14.749495 seconds, including source/key/archive hashing, clones and
negative controls. These phases are nested. Earlier backend compilation, the
original 7.810198-second development key setup and historical controller
execution are separately scoped costs. A later audit adds a 0.006085-second
native verification. This is one shared-host observation with existing keys,
not a replicated speed comparison or evidence that proving beats direct hashing.

All seven frozen controls reject: missing/corrupt receipt bytes, incomplete
leaf inventory, stale parent, reordered events, changed event source commitment
and changed full root. Four exercise the native reader; the last three reject
at exact statement binding before pairing. The proof JSON is 1,484 bytes;
the original 69-file store contains 2,010,722 bytes, including additional
unselected records. A small proof does not replace the original source/receipt
custody that the consumer still reads. The separate audit passes 218 checks,
recomputes all event digests and the full Merkle root, and archives 34 observed
source files. The audit shares the cryptographic backend. This establishes
no new behavioral check, model call, execution reuse, token saving or
production admission.


\paragraph{Actual Requests accepted-source archive.}
A separate post-cohort consumer uses the actual Requests cold-v2 D5414
checkpoint. D6028 failed and contributes no accepted seal. This single event
has four required units: accepted source, native behavior, semantic preparation
and complete non-sealing effects. Two fresh processes read exact archive clones
through native verified-byte storage and admission, recheck original source,
patch, objective, environment, context and effect bindings, and restore all
97 files with exact bytes and raw modes. The original 118-file archive and
ledger remain unchanged. This explicit archive load dispatches no new grader;
the historical controller's ordinary read included a grader replay.

The event digest includes the complete pack, required receipt CIDs and source
inventory. One honest standard v3 proof commits its root/count with three padded
leaves. Proving takes 11.978556 seconds and runtime verification 5.255706
milliseconds; direct producer/consumer read-and-restore take 0.419219 and
0.423702 seconds. The complete 14.737743-second driver includes these costs,
cloning, hashing and all controls. An additional 0.278576-second audit includes
one 7.364014-millisecond verification; previously metered key setup and cohort
execution are separate. The proof is 1,484 bytes, while its supporting archive
occupies 12,514,534 apparent bytes. This consumer retains direct evidence checks;
the proof adds overhead and establishes no amortized or token saving.

Eight frozen mutations reject. Missing/corrupt behavior receipts, omitted units,
changed parent and changed source binding fail archive admission. Changed event,
full root and count fail exact outer-statement binding before pairing. The
changed-parent checkpoint is looked up under delta kind and rejected before
cycle traversal; the initial auditor expected a later kind/cycle error. Its
failure is retained, and a separate corrected auditor checks that exact earlier
lookup rejection without rerunning the experiment. The audit passes 1,462
checks and archives 514 source files. It shares the cryptographic backend.
This is actual repository archive consumption, not inline proof use in the
original cohort, a proof of arbitrary test execution or production admission.

\section{Signed native procedure execution}
\label{sec:signed-native-procedure}
The surviving saved bounded-sequence candidate now passes native seven-layer
verification, certificate issuance, certificate admission and TEST interpreter
execution. Its exact procedure identity is retained. A local TEST-only HMAC
signer and trusted immutable-input custody provide this development admission;
the selection receipt grants no native registry promotion or production
authority. The arithmetic candidate is not relabelled as a repository-repair
procedure.

The prospectively frozen corpus contains 32 test, 24 adversarial, 48 held-out
and 32 shadow inputs, disjoint from the original exposed inputs and one another.
All 136 native executions match independent decisions: 13 true and 123 false.
False is a computed answer, not an execution failure. Six handlers per input
produce 816 dispatches, 408 distinct current check receipts and 136 distinct
decision receipts. The audit reconstructs exact invocation CIDs and restores
all native trace, outcome and receipt objects. These newly authored partitions
were prepared after development; they are not population-generalization evidence.

Eight distinct proof, test, adversarial, held-out, shadow, specification,
source-episode and counterexample artifacts precede signing. Seven bounded Z3
contracts are connected to 4,100 concrete catalog executions; this is a bounded
logical result with a measured source bridge, not a theorem for arbitrary Python.
The original IR has an empty required-proof-contract list. An added adapter
enforces the named proof payload; native contract-specific proof coverage is
not inferred. Independent artifact admission recomputes concrete decisions,
and runtime validation repeats catalog operations. These costs are charged;
this experiment does not demonstrate receipt-only computation avoidance.

Thirty pre-issuer evidence controls and thirteen native runtime controls pass,
including signature, expiry, catalog, observation, query identity, cross-invocation
and domain mutations. LIVE dispatch rejects. Exact invocation checkpoint reuse
avoids redispatch in the same interpreter process; cross-process checkpoint
restoration is untested. Since native canonical values reject floats, a frozen
canonical-JSON transport preserves input bytes behind a checked typed query CID.
The evidence corpus retains 1,192 float-bearing adversaries.

Preparation and freeze take 1.402121 seconds. Three fresh bounded workers take
1.521004 seconds for evidence production, 6.902933 seconds for admission and
issuance, and 8.819003 seconds for execution and controls: 17.242940 seconds
total. The 6.771036 seconds of native invocation calls and 5.127014 seconds of
evidence-negative work are nested. Independent audit adds 2.929698 seconds.
Evidence payloads occupy 2,373,459 bytes and native runtime objects 2,704,141
bytes, excluding snapshots and logs. Earlier float, immutable-mapping and SMT
serialization failures retain their original profiles, limits and costs.
Profile copying before the driver was not timed. The freeze pins 10,736 paths;
156 observed loaded files include 112 standard-library/other paths outside the
native source freeze. Local custody does not supply OS isolation against hostile
Python. No provider call, model-token saving, repository transfer or production
authority is attributed to this integration.
